\begin{papercard}{mulder2019}
\cf{Stages} \stagetag{C} (marketing statements labelled and compared with the policy, the comparison behind RQ3 and RQ4), \stagetag{D} (company websites and policies as separate surfaces), \stagetag{S} (GDPR Art.\ 9 special category and the consent question for health apps). Legal analysis in the European Journal of Law and Technology (Groningen), 2019. Added from the gap search of 23 October 2026. \textbf{Note:} the PDF's text layer drops the letters ``ti'' (``marketng'', ``Artcle''), so passages are quoted as the layer has them.

\cf{Aim} Compare health-app providers' marketing statements about privacy with their policies, and both with the GDPR \pc{mulder2019}{2}{This paper analyses the marketng statements of app providers and the privacy policies of the apps in order to determine whether they are in line with each other and with the GDPR}, because more companies use privacy in their marketing \pc{mulder2019}{3}{a growing number of companies use privacy in their marketng statements}.

\cf{Sample and procedure} Apps named by physicians in a short questionnaire: 34 named, 31 analysed \pc{mulder2019}{3}{In total, 34 diferent apps were mentoned by at least one physician, which were as such selected for this research}. Marketing statements were taken from the companies' public websites \pc{mulder2019}{5}{the public websites of the app companies where investgated to see if they contained any general remarks relatng to privacy}, labelled, and compared with the policies \pc{mulder2019}{3}{I will label the marketng statements of companies with regard to privacy, compare these marketng statements to their privacy policies}.

\cf{Findings} Three groups among general apps: nine with privacy marketing, seven that open their policy with a marketing-like sentence, seven with neither \pc{mulder2019}{6}{Most of the apps - nine in total (see table 1) - mentoned privacy in their marketng statements, demonstratng a positve attude towards privacy}. The policy opener is read as ``a marketing statement in disguise'' \pc{mulder2019}{6}{One could see this as a marketng statement in disguise, especially since all seven companies use such sentences in the beginning of their respectve privacy policy}. Purposes were vague and never linked to the data collected \pc{mulder2019}{11}{Not one policy was clear about the correlaton of the collected data and the purposes for which they are collected}. Conclusion: the marketing promise is not reflected in the policies \pc{mulder2019}{17}{While marketng statements lead you to believe that ‘your’ privacy is important, this is not refected in companies’ privacy policies}; a table linking data, purposes and legal bases is proposed.

\cf{Limits (card)} A single-author legal reading with no coding scheme or reliability; physician-named apps (including general apps such as messengers), not a market sample; policies of spring 2018.

\cf{Use for this thesis} (1) EU precedent for RQ3 and RQ4: privacy claims in marketing and in policy openers are coded as a surface of their own. (2) The policy opener ``your privacy is important to us'' is an Axis C \emph{reassurance} example with a published label (``marketing statement in disguise''). (3) Supports including company websites as a surface in the core corpus, as the methodology does. (4) The missing data--purpose link is a component-level gap Axis A can record.
\end{papercard}
